\documentclass[11pt,a4paper]{article}
\usepackage[T1]{fontenc}
\usepackage[utf8]{inputenc}
\usepackage{lmodern}
\usepackage[margin=25mm]{geometry}
\usepackage{microtype}
\usepackage{amsmath,amssymb,amsthm,mathtools,mathrsfs}
\usepackage{booktabs,array,longtable}
\usepackage{adjustbox}
\usepackage{url}
% File paths may break across lines instead of running into the margin.
\newcommand{\file}[1]{\path{#1}}
\usepackage{enumitem}
\usepackage{xcolor}
\usepackage{graphicx}
\usepackage{pgfplots}
\pgfplotsset{compat=1.18}
\usepgfplotslibrary{groupplots}
\usepackage[round,authoryear]{natbib}
\usepackage[colorlinks=true,linkcolor=blue!50!black,citecolor=blue!50!black,urlcolor=blue!60!black]{hyperref}
\usepackage[capitalise,noabbrev]{cleveref}
% Shared mathematical notation. Arc i -> j means: i requires j.
% Siamo nel mondo dei grafi e delle closure: V sono i vertici, non "item".
\newcommand{\Vset}{\mathcal V}
\newcommand{\Aset}{\mathcal A}
\newcommand{\Gset}{\mathcal G}
\newcommand{\Cset}{\mathscr C}
\newcommand{\Rset}{\mathbb R}
\newcommand{\Sseq}{\mathcal S}
\newcommand{\Fset}{\mathcal F}
\newcommand{\Qset}{\mathcal Q}
\newcommand{\Bset}{\mathcal B}
\newcommand{\Nset}{\mathcal N}
\newcommand{\Tset}{\mathcal T}
\newcommand{\Kset}{\mathcal K}   % blocchi canonici K_1, ..., K_k
% Matrici e vettori in grassetto; le lettere greche via \boldsymbol.
\newcommand{\mat}[1]{\mathbf{#1}}
\newcommand{\vect}[1]{\mathbf{#1}}
\newcommand{\vecg}[1]{\boldsymbol{#1}}
\newcommand{\ind}{\chi}
\newcommand{\divg}{\operatorname{div}}
\newcommand{\supp}{\operatorname{supp}}
% Nome del metodo (decisione D18): minuscolo, sul modello di *network simplex*.
% \FS in mezzo alla frase, \FSC a inizio frase, \CS per la forma breve.
\newcommand{\FS}{ratio-closure simplex}
\newcommand{\FSC}{Ratio-closure simplex}
\newcommand{\CS}{closure simplex}
\newcommand{\LP}{\textsc{LP}}
\newcommand{\rhoT}{\rho}


% Generated artefacts are produced by tools/report/analyze.py from the stored
% campaigns.  A campaign that has not been run yet must not break the build:
% the document then states which artefact is missing instead of showing numbers.
\newcommand{\generated}[1]{%
  \IfFileExists{generated/#1.tex}{\input{generated/#1}}%
  {\begin{center}\small\itshape Artefact \texttt{\detokenize{generated/#1.tex}} not present:
   the corresponding campaign has not been analysed in this build of the report.\end{center}}}
\newcommand{\generateddata}[1]{generated/#1.dat}

\pgfplotsset{
  every axis/.append style={
    width=0.46\textwidth, height=0.34\textwidth,
    tick label style={font=\footnotesize}, label style={font=\small},
    legend style={font=\footnotesize, draw=none, fill=none},
    grid=both, grid style={line width=0.2pt, draw=gray!25},
  },
}

\title{A ratio-closure simplex for maximum-ratio closure\\ and its canonical decomposition:\\ implementation and experimental report}
\author{Experimental report generated from the stored campaigns}
\date{\today}

\begin{document}
\maketitle



\begin{abstract}
This report documents the measurement campaigns run on the \FS{} solver of this
project and on the two independent combinatorial backends built as references.
It states the questions asked, the instances and protocol used, the validation
performed on every result, how the solver's parameters were tuned on data
disjoint from the test set, and what the measurements do and do not support.
Every table and figure is regenerated from the stored run manifests by
\texttt{tools/report/analyze.py}. The figures quoted in the running text are
taken from those same tables and from the stored manifests, and were re-checked
against them after the instance library was regenerated deterministically and
every campaign re-run (\file{docs/decisions.md}, D16).
\end{abstract}

\tableofcontents

\section*{Summary of findings}
\addcontentsline{toc}{section}{Summary of findings}

\paragraph{What was measured.} One linear program, three independent solvers.
The \FS{} of this project, a Dinkelbach iteration whose ratio oracle is a
minimum cut, and a divide-and-conquer parametric decomposition. All three return
the same object --- an optimal LP solution, its multipliers and the canonical
sequence of blocks --- and every result is re-checked by a separate process
before it is counted.

\paragraph{Correctness.} The three backends agree everywhere they were compared,
and agree with an exhaustive rational oracle on small instances and with an
independent LP solver on larger ones. Nothing in the campaigns contradicts the
correctness of any backend.

\generated{table_summary}

\paragraph{Speed.} The parametric backend is the only configuration that solves
the whole test set: $46$ of $46$. Dinkelbach solves $40$, the tuned \FS{} $36$
and the reference \FS{} $27$. Two families are solved by the parametric backend
alone --- disconnected graphs and the instances derived from MineLib --- where
the other three finish nothing within the limit.

The medians do not rescue the counts. Over the instances it solves, the tuned
\FS{} has a median of $0.142$\,s against the parametric backend's $0.057$\,s,
and that median is taken over the smaller subset it can finish; Dinkelbach sits
between them at $0.234$\,s but solves four more instances. On the historical
precedence-knapsack instances all four configurations except the reference
finish everything, and the ordering there is unambiguous: $0.027$\,s median for
the parametric backend, $0.087$\,s for Dinkelbach, $0.655$\,s for the tuned
\FS{}. The honest summary is that this implementation is neither the fastest
where it finishes nor able to finish everywhere.

\paragraph{What made the difference.} Not the pricing rule. The first
measurements had \FS{} one to three orders of magnitude behind, and parameter
search over pricing schemes moved the result by a few per cent. What moved it by
factors was the handling of degeneracy: the primal spends a median of $99.3\%$
of its pivots on zero steps, and three changes follow from that ---
releasing the anti-cycling fallback instead of letting it become the operating
rule, restricting the ratio-test scan to the arcs that can actually block, and
stopping that scan at the first zero-step blocker. The last two reproduce
techniques already present in the group's historical prototype.

The ablation confirms the ranking and corrects two of the search's choices.
Measured as paired speedups against the setting that forces Bland's rule,
the frozen configuration is $2.83$ times faster at the median and up to $18.8$
times on individual instances, so the anti-cycling release is worth a factor
rather than a percentage; moving the degeneracy trigger back to its old value
gives up most of that, $1.25$ instead of $2.83$. The ratio-test refinements are
worth about a factor of two each, $2.17$ and $2.20$.

Two settings the search did \emph{not} choose are faster than the one it did:
starting from the sink basis scores $3.33$ and enabling transitive reduction
$3.46$, against the frozen configuration's $2.83$. The frozen configuration is
reported as it was frozen, on training and validation data only, and not
silently improved after the fact --- but the ablation says plainly that the
search left something on the table, and \cref{sec:discussion} treats it as a
finding rather than a footnote.

\paragraph{How to read the tables.} Times are wall-clock medians over replicas of
the same instance and configuration. \emph{PAR-2} charges twice the time limit
for every instance a configuration did not solve, so it mixes speed and
robustness into one number; it is always reported next to the plain count of
solved instances so that a time-out is never hidden inside an average. Paired
speedups are computed only on the instances both settings solved and verified.
Performance profiles state the compared set and never count an unsolved instance
as covered.

\paragraph{What this report does not claim.} It does not claim novelty: the
audit of the literature was deliberately not carried out in this phase. It does
not claim anything about the state of the art in maximum flow: the flow backends
use a Dinic implementation written for this project, not a competitive flow code.

\section{Objectives and hypotheses}
\label{sec:objectives}

The object of the evaluation is the \FS{} solver: a combinatorial simplex whose
basis is represented by a forest on the precedence graph, applied to the
normalised maximum-ratio subproblem and iterated on the residual to build the
canonical sequence of blocks that yields the LP value at any capacity.

Two further backends were implemented for this report and are used as
references, not as contributions: a Dinkelbach iteration whose ratio oracle is a
maximum-closure minimum cut, and a divide-and-conquer parametric decomposition
of the same closure family. All three produce the same certified structures and
are checked by the same verifiers.

The campaigns answer the following questions.

\begin{description}[leftmargin=2.4em,style=nextline]
  \item[E1 Correctness] Does the solver produce certified optimal solutions on
    generic DAGs, with negative profits, ties and degeneracy? Do the three
    backends and an independent LP solver agree?
  \item[E2 Overall performance] How do the reference \FS{} configuration, the
    tuned one and the two combinatorial baselines compare on a frozen test set,
    per family rather than only in aggregate?
  \item[E3 Scalability] How do time and memory grow with $n$ and $m$, and how
    much of the growth is the cost of a pivot rather than the number of pivots?
  \item[E4 Ablation] What does each implementation choice contribute: local
    forest updates, dual canonicalisation, partial pricing, transitive
    reduction, vertex ordering?
  \item[E5 Tuning] Does the search produce a default that transfers from
    training and validation to instances never used to choose it?
  \item[E6 One capacity against many] When does building the whole positive path
    pay for itself against a single solve, once the cost of the queries and of
    materialising the solution is included?
  \item[E7 Degeneracy and numerics] What is the cost of ties, exact and
    near-exact, and of the anti-cycling fallback?
  \item[E8 Applied instances] What happens on the historical PCKP instances and
    on the open-pit instances derived from MineLib?
\end{description}

No question is posed so that a particular answer is the only admissible one. In
particular, the report states plainly where the \FS{} implementation is slower
than the baselines, and section~\ref{sec:discussion} discusses why.

\section{Methods}
\label{sec:methods}

\subsection{Problem and required output}

All backends solve the same linear program: maximise $\sum_i p_i x_i$ subject to
$\sum_i w_i x_i \le c$, $x_i \le x_j$ for every arc $(i,j)$ meaning ``$i$
requires $j$'', and $0 \le x_i \le 1$. Profits carry either sign, weights are
strictly positive, and there is a single capacity row. Cyclic digraphs are
accepted and condensed; solutions and multipliers are expanded back to the
original variables and arcs.

The required output of a \texttt{solve} task is a primal solution, the
multipliers $\lambda$, $\mu$ and $\alpha$, and the canonical blocks of the
positive path. Comparisons are therefore between runs that produce the same
object: whenever a configuration answers several capacities from one path, the
report says so explicitly and section~\ref{sec:scalability} treats it separately.

\subsection{Compared methods}

A note on the labels. The method developed by this project is the \FS{}, and
the configuration labels in the current suites are named after it,
\texttt{ratio-closure-*}. The labels below instead read \texttt{forest-*},
because that is what the campaigns reported here actually recorded: they were
measured on 6 September 2026, before decision~D18 was applied to the code, and a
run identifier is derived from the instance, the task, the capacity, the
configuration \emph{content} and the hash of the sources. The manifests are a
record of what was executed and are not rewritten, so this report quotes the
labels as measured. The divergence is cosmetic and disappears with the first
campaign measured under the current suites; the analysis script recognises both
spellings.

\begin{description}[leftmargin=2.4em,style=nextline]
  \item[\texttt{forest-reference}] \FS{} with full pricing, Bland's entering
    rule, complete forest rebuilds and sequential canonicalisation. This is the
    verifiable reference implementation, not a configuration meant to be fast.
  \item[\texttt{forest-tuned}] the configuration frozen by the random search of
    section~\ref{sec:tuning}, chosen on training and validation data only.
  \item[\texttt{dinkelbach}] Newton iteration on
    $\max\{p(S)-\lambda w(S)\}$ over closures, started from a strict lower bound
    on the ratio so that negative ratios need no special case. Each iterate is a
    maximum-closure minimum cut.
  \item[\texttt{parametric}] divide and conquer on the interval of ratios: the
    mean ratio of an interval is either a breakpoint or splits it in two. Each
    subproblem is an independent maximum flow.
\end{description}

The maximum-flow routine is a Dinic implementation written for this project
(level graphs, current arc, iterative augmentation). It is correct and
reasonable, but it is \emph{not} a state-of-the-art flow code: no pseudoflow, no
highest-label push-relabel with gap and global relabelling, no reuse of flow
across parametric subproblems. Any statement about the flow backends in this
report is a statement about this implementation, and
section~\ref{sec:discussion} says what that does and does not license.

\subsection{Accuracy is fixed before speed}

All configurations use the same tolerances, and tolerances are not part of the
tuning space. Every run must pass, in the original units, the primal checks
($0 \le x \le 1$, every original arc, the capacity row), the dual sign
conditions, the reduced-cost inequality
$w_i\lambda + \mu_i + \divg_i(\alpha) \ge p_i$, complementarity on $\mu$, on
$\alpha$ and on the capacity row, and equality between the primal objective and
the dual bound. A result that fails any of these is reported as
\texttt{numerical\_failure} and never counted as a success.

Verification is a separate process: after each run, \texttt{pclp verify} reloads
the instance and the saved result and re-checks the certificate without
re-optimising. Only runs that pass this second, independent check enter the
analysis.

\section{Environment}
\label{sec:environment}

\generated{table_environment}

All measurements come from a single machine and a single build. The solver is
compiled in \texttt{release} mode without link-time optimisation and without
\texttt{-march=native}, so that the reported times refer to a portable build;
sanitiser builds and per-pivot tracing exist but are never used for
measurement. Each solver process is pinned to one core, and campaigns are run one at a time.
This is not a precaution taken on principle: during the preliminary exploration
the same configuration on the same instance measured about $40$ seconds with a
single active process and over $60$ seconds with three processes pinned to
distinct cores of the same machine. Pinning does not remove contention for
shared cache and memory bandwidth, so concurrent campaigns were ruled out and
the wall-clock cost was accepted instead.

The build fingerprint stored in every result contains the SHA-256 of all sources
that enter the binary, the compiler, the build type and the flags. A run whose
build hash differs from the campaign's is not counted in that campaign: changing
the code produces new run identifiers rather than silently mixing measurements.

Wall time is measured by the runner with a monotonic clock around the whole
process, including reading the instance and writing the result. The solver
separately reports the time of its own phases. Peak memory is the solver
process's \texttt{ru\_maxrss}, kept distinct from the runner's memory.

\section{Instances}
\label{sec:instances}

\generated{table_instances}

\subsection{Generated families}

Each synthetic family isolates one structural phenomenon: absence of arcs
(overhead only), chains (maximum depth), precedence forests, layered DAGs with
and without level-skipping arcs, sparse and dense random DAGs, diamonds and
shared prerequisites, disconnected components, transitively redundant arcs,
artificial strongly connected components, and two-level graphs. Across these,
the generator varies the size, the weight spread (uniform, dispersed,
heterogeneous by orders of magnitude), the share of negative profits, the
correlation between profit and weight, and the share of exact ties. Every
instance stores its generator, parameters and seed, so the whole set is
rebuilt byte for byte by one command.

\subsection{Historical PCKP instances}

Twenty-three single-capacity precedence-constrained knapsack instances were
imported from the project's historical archive, with $349$ to $11\,757$ vertices and
up to $83\,218$ arcs. Each was validated against the companion CPLEX \texttt{.lp}
model: the importer accepts only one capacity row plus simple precedence rows
with unit bounds, and refuses everything else. Because the two files do not
always share the variable numbering, the check falls back to colour refinement
and claims a correspondence only when every colour class is a singleton and the
mapped arc sets coincide exactly. The multi-period files of the same archive
have one capacity row per period and are therefore outside this family; they
were refused, not adapted.

\subsection{Open-pit instances derived from MineLib}

Eight open-pit instances, up to $99\,014$ blocks and $1\,271\,207$ precedence
arcs, were derived from the MineLib CPIT models. The derivation is explicit and
recorded inside every imported instance: period indices and discounting are
dropped, the knapsack row is resource~0, and the capacity is that resource's
limit summed over the periods it was granted. These are realistic structural
benchmarks, not the original MineLib models, and the report never presents them
as such. Three MineLib instances were refused because some blocks have no
resource-0 coefficient and would carry zero weight, outside the supported family.

\subsection{Frozen splits}

Training, validation and test manifests were written once, stratified by origin,
family and size decade, keeping all replicas of one generated base on the same
side. Parameter search uses training only, the finalists are compared on
validation, the chosen configuration is frozen with its date and rationale, and
only then is the test set measured.

\section{Validation}
\label{sec:validation}

\generated{table_validation}

Correctness is established at three levels, each with a calculation independent
of the solver being checked.

\paragraph{Exhaustive rational oracle.} For small instances the acceptance suite
enumerates every subset, keeps the closures, and computes with exact fractions
the maximum ratio, \emph{all} its maximisers and their canonical union, then
repeats on the residual to obtain the whole sequence. The LP value is obtained
independently from the convex hull of the points $(w(C),p(C))$, including pairs
of closures straddling the capacity, without assuming that the two are nested.
The solver must reproduce the same blocks, the same ratios and the same LP
values.

\paragraph{Basis algebra in exact arithmetic.} For traced runs the suite
rebuilds, in exact arithmetic, the explicit basis matrix of every pivot and
recomputes the basic values, the direction, the reduced cost, the ratio-test
step, the leaving variable and the pivot coefficient. It requires that all four
structural pivot types --- vertex or slack entering against vertex or slack leaving
--- actually occur.

\paragraph{Independent LP solver.} The same instances are handed to HiGHS
through SciPy as a plain LP with one variable per vertex and one row per arc. Only
the optimal value is compared: the optimal vector need not be unique.

\paragraph{Cross-backend agreement.} The Dinkelbach and parametric backends
produce their own certificates from their own flows, and are required to return
the same canonical blocks and the same LP values as \FS{}. Because they share
no code with the pivot loop, an error in the forest engine cannot be masked by an
identical error in its verifier.

\paragraph{Metamorphic checks.} Permuting vertex identifiers and arc order,
duplicating precedence rows, adding transitively redundant arcs, adding an
isolated negative vertex, scaling profits, and scaling weights and capacity
together all leave the value and the canonical blocks where they should be.
Condensing an artificial strongly connected component must agree with the
equivalent aggregated vertex.

\paragraph{Per-run verification.} Beyond the suite, every run of every campaign
is re-checked by a separate \texttt{pclp verify} process that reloads the
instance and the saved result and validates the certificate without
re-optimising. The analysis counts only runs that pass this second check.

\section{Tuning}
\label{sec:tuning}

\subsection{Search space and protocol}

Only parameters the solver actually implements take part in the search:
pricing scheme and its block size or candidate limit, entering rule, forest
update mode with its rebuild interval and fraction, canonicalisation,
degeneracy trigger, vertex ordering of the condensation, transitive reduction and
the initial basis. Conditional parameters are sampled only when their guard
holds, so a candidate never carries a block size it cannot use, and the
constraint that Bland's rule requires full pricing mirrors the solver's own
validation. Numerical tolerances are deliberately excluded: relaxing them would
buy time by solving a different problem.

Candidates are drawn by reproducible random search from a fixed seed, after a
small list of a-priori plausible configurations. Each candidate is evaluated on
a deterministic subsample of the training manifest under a fixed per-run limit,
and ranked lexicographically by number of incorrect results, number of instances
solved and verified within the limits, PAR-2 time, and peak memory. Every
attempt, including rejected ones and the reason for rejection, is appended to a
history file, so an interrupted search resumes without repeating work and the
total cost of tuning --- not only the cost of the winning configuration --- is
recorded.

\generated{table_tuning_forest-v2}

A first search was run before the parameters that turn out to matter existed in
the solver, and its attempts differed by less than three per cent: it is recorded
in \file{experiments/tuning/} but it selected nothing and is not reported here.
The search shown above is the one whose winner was frozen.

\subsection{Frozen configuration}

The configuration selected on training and validation data is written to
\file{configs/frozen/} together with its date, the split it was chosen on, the
search space and the rationale, and only then is the test manifest measured. The
test set is never used to adjust a default.

\section{Main results on the frozen test set}
\label{sec:results}

\generated{table_main}

\generated{table_main_speedup}

\generated{table_agreement}

The comparison is between runs that produce the same object: a certified primal
solution, the multipliers, and the canonical blocks of the positive path at
the instance's own capacity. Instances that no configuration solved within the
limit are still counted in PAR-2 and in the solved column; they are not removed
from the sample.

\begin{figure}[tbp]
  \centering
  \IfFileExists{generated/profile_main.dat}{%
    \begin{tikzpicture}
      \begin{axis}[xlabel={$\tau$ (factor over the best time on the instance)},
                   ylabel={fraction of instances}, xmode=log, ymin=0, ymax=1.02,
                   legend pos=south east, width=0.62\textwidth, height=0.42\textwidth]
        \addplot table[x=tau, y index=1] {generated/profile_main.dat};
        \addplot table[x=tau, y index=2] {generated/profile_main.dat};
        \addplot table[x=tau, y index=3] {generated/profile_main.dat};
        \addplot table[x=tau, y index=4] {generated/profile_main.dat};
        \legend{\texttt{dinkelbach},\texttt{forest-reference},\texttt{forest-tuned},\texttt{parametric}}
      \end{axis}
    \end{tikzpicture}}{\small\itshape Performance profile not available: the main
    campaign has not been analysed in this build of the report.}
  \caption{Performance profile on the frozen test set. A configuration is
  counted at $\tau$ when its median time on an instance is within a factor
  $\tau$ of the best median time recorded on that instance; unsolved instances
  never enter the numerator. The compared set is the instances at least one
  configuration solved and verified.}
  \label{fig:profile-main}
\end{figure}

\subsection{Against the 2023 prototype}
\label{sec:legacy}

The group's 2023 \FS{} prototype was imported into this project under
\file{legacy/forest2023/} and measured on the same maximum-ratio oracle, which is
the task it implements. Three patches were needed to run it at all and are
recorded in that directory: an absolute CPLEX include replaced by opaque
typedefs, four interactive \texttt{cin.get()} pauses removed, and the append to a
file in the working directory removed. No line of the algorithm was changed.

Running it exposed limits that belong to the prototype, not to the comparison.
It stores profits and weights as \texttt{int}, so instances with fractional data
cannot be represented and are refused instead of rounded. It allocates two
$n\times n$ adjacency matrices, so its memory is quadratic in the number of
vertices. Its initial basis assumes a spanning tree of $n-1$ edges, so it requires a
connected graph. Finally, the feasibility shortcut added in the 2023-05-11
version recurses without termination on a generic, non-bipartite DAG --- it
overflows the stack already on the eight-node example --- so the measurements use
it disabled, which reproduces the behaviour of the 2023-05-05 version of the same
file.

\paragraph{The outcome.} On the eighteen instances measured, the prototype
completed \emph{none}, so there is no time to compare it on and no table worth
printing. Four times it hit the time limit --- including a chain of two hundred
vertices, which the current solver certifies in nine milliseconds --- once it
crashed, once it exceeded the memory budget, and the remaining twelve it refused
outright, either because the data are not integral or because the graph is
disconnected and no spanning tree of $n-1$ edges exists. The current solver
returned a certified optimal ratio on every one of the eighteen, in times between
seven milliseconds and thirty-two seconds. The per-instance record is in
\file{experiments/analysis/legacy\_baseline*.json}.

That is a wide gap, and it should be read for what it is. Most of it is not a
difference in pivot rules: it is that the prototype rebuilds the whole forest and
reallocates two arrays at every iteration, recurses instead of iterating, and
carries data structures quadratic in the number of vertices. Those are the costs of
a research prototype that was never asked to run outside the bipartite instances
it was written for, and the current solver inherits none of them. The comparison
therefore measures engineering, not the merit of the underlying method --- which
is the same method in both.

What the prototype did contribute is two ideas that survived measurement and are
in the current solver: the early exit of the ratio test on a zero-step blocker,
which it implemented as an explicit feasibility check during the forest
traversal, and the absence of a permanently engaged anti-cycling rule, which the
ablation of section~\ref{sec:analysis} confirms was costing a factor in pivots. A
third idea is still on the table and not yet implemented: the Euler-tour
intervals it maintained, which give subtree membership in constant time and would
attack the rebuild cost that now dominates.

Where both solvers return an optimal ratio, they return the same one. On this
instance set that condition never held, so the cross-check rests on the
eight-node example of the tutorial, where both return the ratio $2$.


\section{The full canonical decomposition}\label{sec:canonical}

Every campaign in \cref{sec:results} answers a single capacity: the solver walks
the canonical sequence only as far as the capacity requires. This section asks
the other question --- produce the \emph{entire} sequence of closure layers,
task \texttt{full-path} --- because that is the object the method is designed to
deliver, and because it is the setting in which a basis, unlike a flow state,
could in principle be carried from one residual to the next.

The suite is \file{configs/suites/canonical\_path.json}: twelve instances across
the generated families, one historical PCKP instance and one MineLib instance,
two replicas, a sixty-second limit.

The same canonical sequence, restricted to precedence graphs that are directed
forests, is the subject of the companion work \citep{DoseEtAl2026}, whose
aggregation algorithms reach $\mathcal{O}(n\log n)$ and whose implementation is
public at \url{https://github.com/fabiofurini/parametric-closure-forests}. What
this report calls a canonical block is a \emph{closure layer} there, and the
per-vertex bound computed here is their \emph{threshold} of a vertex. The
comparison below is on general DAGs, where those algorithms do not apply.

\generated{table_canonical_path}

\generated{table_canonical_path_speedup}

\paragraph{Restarting each oracle cold.} With every residual oracle rebuilt
from the empty basis --- configurations \texttt{ratio-closure-tuned} and
\texttt{ratio-closure-dual} --- the method is not competitive. The parametric
backend is faster on every instance but the arc-free graph, the median ratio is
a factor of twelve, and on a layered DAG with $n = 2000$ it is a factor of
$859$. The oracle count is exactly what the design predicts, one per closure
layer, and the parametric backend uses about two flow computations per layer.
The difference is inside the oracle: each one costs tens to hundreds of
thousands of pivots, of which between $98\%$ and $100\%$ are degenerate.

\paragraph{Restarting each oracle from the previous final basis.} The extracted
block \emph{is} the positive component, that is one tree of the basis forest;
the other trees survive untouched and stay anchored at zero. Reusing that forest
is what \texttt{--warm-start residual} does, and the correctness condition is a
single test: since the values are $1/w(T)$ on the positive tree $T$ and zero
elsewhere, an active arc leaving $T$ would violate feasibility, so $T$ must be
closed in the residual. The best closed surviving tree is chosen; if none is
closed the oracle falls back to a cold restart. The check costs $\mathcal{O}(m)$
per oracle and is verified at every pivot in the acceptance suite.

The effect is on the pivot count, which is where the time was:

\begin{center}
\begin{tabular}{@{}lrrr@{}}
\toprule
instance & cold pivots & warm pivots & factor \\
\midrule
\texttt{gen\_disconnected\_components6\_n2000} & 251{,}780 & 3{,}080 & $81.7\times$ \\
\texttt{gen\_sparse\_degree3\_n2000}           &  89{,}468 & 6{,}108 & $14.6\times$ \\
\texttt{gen\_transitive\_degree2-extra3.0\_n2000} & 166{,}730 & 13{,}225 & $12.6\times$ \\
\texttt{pckp\_A\_972\_1661}                    &   8{,}503 & 1{,}056 & $8.1\times$ \\
\texttt{minelib\_newman1}                      &  48{,}974 & 6{,}915 & $7.1\times$ \\
\bottomrule
\end{tabular}
\end{center}

The decomposition produced is identical block by block, and the objective of
\texttt{solve} is unchanged to nine decimals.

\paragraph{The ranking, with the warm restart in.} On the ten instances every
configuration solves within the limit, the warm primal is $3.1$ times faster
than the cold one in total time and overtakes Dinkelbach, which the cold
configuration did not. The simplex now wins three of the ten instances --- the
chain and the disconnected families, where the sequence has many layers, and the
cyclic one --- against one out of twelve before. The parametric backend still
wins the remaining seven and is $3.4$ times faster in total, so the conclusion
of this report is unchanged in direction and changed in size: the gap is no
longer one to two orders of magnitude but a factor of three, and it is closed
entirely on the families with the longest sequences.

\paragraph{What this does and does not settle.} It settles that the cold
sequential decomposition was leaving most of its work on the table, and that the
structure of the method --- a basis, not a flow --- is what recovers it. It does
not settle the comparison against parametric flow, which still wins on the
majority of instances and in total time. Two things are now worth measuring and
are not measured here: the frozen configuration was tuned on \texttt{solve} and
never on \texttt{full-path}, and a paired ablation on this task puts
\texttt{initial\_basis=sink} $1.21$ times ahead of it; and the degeneracy that
dominated the cold oracles is reduced, not removed, so the anti-cycling regime
deserves its own tuning on this task.

\section{Where the time goes}
\label{sec:analysis}

\generated{table_main_breakdown}

\generated{table_ablation}

\generated{table_ablation_speedup}

The breakdown separates the cost of one pivot from the number of pivots. The
first is bounded by the structures: full pricing scans the $n-1$ non-basic
variables of the residual, the direction visits the entering subtree and the
positive component, and the ratio test examines every basic slack of the
residual, so a pivot costs $O(n+m)$ regardless of how small the change to the
forest is. The second is a property of the instance and of the entering rule,
and it is the quantity the ablation is designed to expose.

The ablation switches off one implementation choice at a time with respect to the
frozen configuration, on the validation manifest, with the same data and the same
acceptance thresholds as the main campaign.

\subsection{What separates the easy families from the hard ones}

It is tempting to attribute the cost to the number of blocks, since each one
costs a fresh oracle with a fresh basis. The measurements do not support that.
Across generated instances of two thousand vertices, all with the frozen
configuration: a chain produces $11$ blocks and $4\,640$ pivots in $0.116$\,s;
a sparse DAG of out-degree three produces $282$ blocks and $134\,830$ pivots in
$1.311$\,s; a layered DAG of twenty levels produces $232$ blocks and
$348\,205$ pivots in $6.079$\,s; and a layered DAG of eight levels with dense
connections between them produces only $85$ blocks --- fewer than either sparse
instance --- yet $528\,841$ pivots, and does not finish within the limit. More
blocks is not more work.

The quantity that separates them is pivots \emph{per oracle}, expressed as a
multiple of $n$: about $0.21n$ on the chain, $0.24n$ on the sparse DAG,
$0.75n$ on the twenty-level instance and $3.1n$ on the eight-level dense one.
The layered structure gives every entering variable many precedence slacks
between vertices that are all at zero, so each oracle spends thousands of
zero-step pivots walking to its optimum.

Two consequences follow for where effort should go. Reusing the basis across
residual calls, which would save one initialization per oracle, cannot help the
hard case: a hundred initializations against six hundred thousand pivots. What
can help is either making each pivot cheaper --- incremental pricing and updates
confined to the modified paths, since pricing and rebuilding are together four
fifths of the solver's time on this family --- or attacking the number of
degenerate pivots directly with a lexicographic or perturbed ratio test. Only the
second changes the multiplier, and it is the one with no implementation here.

This also explains the shape of the comparison with the parametric backend. That
backend pays nothing for degeneracy: it solves one maximum flow per subproblem on
disjoint vertex sets. Its advantage is largest exactly on the family where the
forest's pivot count per oracle is largest, which is not a coincidence of
constants but a difference in what each method spends its work on.

\subsection{The dual simplex}
\label{sec:dual-results}

\generated{table_dual}
\generated{table_dual_speedup}

The measurements above say that the primal loop is dominated by zero-step
pivots. Primal degeneracy and dual degeneracy are different defects, and here
they are wildly asymmetric: the primal kind is a per-pivot phenomenon, at a
median of $99.3\%$ of pivots over the whole dual campaign, while the only source
of the dual kind is a tie between block ratios, which is a per-block event and
therefore bounded by the number of blocks. That asymmetry is the standard
indication for the dual method, and it was implemented and measured.

Two structural facts make it unusually clean on this problem. The vertex variables
are never primal infeasible, since their value is $1/w(P)$ on the positive
component and zero elsewhere, so every primal infeasibility is a precedence
slack, that is an arc leaving $P$; and the basis is primal feasible exactly when
$P$ is closed, which is therefore the stopping rule. The starting basis --- every
vertex its own tree, $P$ the vertex of maximum ratio --- is dual feasible by
construction and costs $O(n)$.



The prediction is confirmed on the quantity it was made about, and refuted on
every other. Table~\ref{tab:dual} and table~\ref{tab:dual-speedup} report the
dedicated campaign on the frozen validation split, $31$ instances in two
replicas under the same protocol as every other campaign.

Degeneracy collapses: the primal sits at a median of $99.3\%$ of pivots with
zero step, the dual at $0.0\%$. Nothing else improves. The dual performs
\emph{more} pivots, not fewer --- a median of $7\,813$ against the primal's
$3\,253$ on the instances both solve --- and it is slower: the paired speedup
against the primal has a median of $0.90$, ranging from $0.34$ to $2.47$. It
also solves less: $56$ of $62$ runs against $62$ of $62$, losing all three
replicas of a dense DAG of two thousand vertices that the primal solves in
about eleven seconds. The variant that picks the largest tree is worse again,
median $0.78$ and $54$ of $62$.

So the reduction of degeneracy was real and worth nothing. The dual traded null
steps for a longer path, which is the outcome that the primal's own diagnosis
should have predicted: the cost was never the null steps as such, it was the
number of pivots needed to walk to the optimum, and the dual walks farther. On
the same instances the parametric backend is $3.89$ times the primal at the
median and up to $135$ times on one of them.

Choosing the leaving arc by a structural criterion instead of by smallest index
does not help either, and there is a structural reason: every violated slack has
the same value $-1/w(P)$, so a most-infeasible rule has nothing to discriminate.
The largest-tree rule is the one carried into the campaign, and it is the slowest
of the three settings measured.

Correctness is not at stake in any of this. The dual returns the same canonical
blocks as the primal on the whole exhaustive suite and on fifteen hundred
random instances checked with the diagnostics enabled at every pivot, and it
produces its certificate through the same verified code. If dual feasibility is
lost to rounding, the solver restarts from the primal initialization rather than
inheriting an infeasible basis, so the answer is certified either way.


\section{Scalability and many capacities}
\label{sec:scalability}

\generated{table_scalability}

\begin{figure}[tbp]
  \centering
  \IfFileExists{generated/scaling_scalability_parametric.dat}{%
    \begin{tikzpicture}
      \begin{axis}[xlabel={vertices $n$}, ylabel={median seconds}, xmode=log, ymode=log,
                   legend pos=north west, width=0.62\textwidth, height=0.42\textwidth]
        \IfFileExists{generated/scaling_scalability_forest-tuned.dat}{%
          \addplot table[x=nodes, y=seconds] {generated/scaling_scalability_forest-tuned.dat};
          \addlegendentry{\texttt{forest-tuned}}}{}
        \IfFileExists{generated/scaling_scalability_dinkelbach.dat}{%
          \addplot table[x=nodes, y=seconds] {generated/scaling_scalability_dinkelbach.dat};
          \addlegendentry{\texttt{dinkelbach}}}{}
        \addplot table[x=nodes, y=seconds] {generated/scaling_scalability_parametric.dat};
        \addlegendentry{\texttt{parametric}}
      \end{axis}
    \end{tikzpicture}}{\small\itshape Scalability data not available: the
    scalability campaign has not been analysed in this build of the report.}
  \caption{Median wall time against the number of vertices on the generated
  families whose structure is held fixed while the size varies. Only runs that
  finished and were verified appear; a configuration whose curve stops has hit
  the time limit beyond that size.}
  \label{fig:scaling}
\end{figure}

\subsection{One capacity against many}

\generated{table_capacities}

The canonical positive path is built once and answers any capacity: locating the
critical block costs $O(\log k)$ in the number of blocks, but exporting a
dense solution vector costs at least $O(n)$ per query. The campaign therefore
separates the cost of the path from the cost of the queries, and reports the
number of queries at which building the whole path pays for itself against
solving each capacity independently.

\section{Degeneracy, ties and numerics}
\label{sec:robustness}

\subsection{Why the pivot loop stalls}

The normalised subproblem is degenerate by construction. In a basis the vertex
values are constant on the positive component and zero on every other tree, so a
step is blocked at length zero whenever the blocking variable is one of the many
already at zero: a vertex of the entering subtree with a negative direction
component, or a precedence slack between two vertices that are both at zero. On the
applied instances the median share of pivots with zero step is $99.3\%$, and it
reaches $99.9\%$; on the family with no arcs at all, where no slack can block,
the share is zero. Degeneracy is
therefore a property of the precedence structure, not an accident of the
implementation.

This observation drove three changes, each kept as a separate option so that
section~\ref{sec:analysis} can measure them one at a time.

\paragraph{The anti-cycling fallback must not be the operating policy.} Switching
permanently to Bland's rule after a run of degenerate pivots made every entering
rule collapse onto Bland within the first oracle, because the run of degenerate
pivots never ends. Moving the trigger far away, and releasing the fallback at the
first pivot with a positive step, restores the improving rule without giving up
termination: an episode under Bland runs on a fixed basic solution and cannot
cycle, and every episode ends with a strict increase of the ratio, so only
finitely many episodes are possible.

\paragraph{The ratio test need not scan every arc.} The direction is $+\sigma$ on
the entering subtree, $-\kappa$ on the positive component and zero elsewhere, so
an arc can block only if it crosses the boundary of one of the two sets. Because
the positive component is closed, an arc that crosses it has its tail outside,
which makes the restricted scan complete.

\paragraph{A zero step is the minimum, so the scan can stop at the first one.}
This is the shortcut the historical prototype implemented as an explicit
feasibility check during the forest visit. It is suspended while Bland's rule is
active, because that rule needs the smallest index among ties.

\subsection{Ties and numerics}

Exact ties are handled by construction: the canonical block is the union of
all maximisers, obtained either by complementarity from the dual or by
aggregating consecutive increments of equal ratio. Both settings return the same
blocks on the whole exact suite. Near-ties are separated by the relative
optimality tolerance and are deliberately not merged.

All arithmetic is \texttt{long double} with binary scaling, which is exact in
binary floating point and does not turn an integer zero-profit sum into a
spurious nonzero ratio. Certificates are re-checked in the original units, after
expanding strongly connected components, so a numerical success in the scaled
model cannot pass as a success in the user's model.

\section{Discussion}
\label{sec:discussion}

\subsection{What the measurements support}

The three backends agree on every instance of the exact suite and on every
instance of the campaigns they both solved: same canonical blocks, same LP
values, each with its own certificate produced by code that shares nothing with
the others. Together with the independent LP solver and the exhaustive rational
oracle, this supports the correctness of the implementation on the domain it
claims: generic DAGs, profits of either sign, strictly positive weights, one
capacity, cyclic digraphs by condensation.

On performance the picture is uneven, and it is given per family because a
single number would hide the shape of it. The parametric backend solves the whole
test set; Dinkelbach solves $40$ of $46$, the tuned \FS{} $36$ and the reference
\FSC{} $27$. The tuned \FS{} is not the fastest even where it finishes: its
median over solved instances, $0.142$\,s, is above the parametric backend's
$0.057$\,s, and it is taken over the smaller set it can finish.

Per family, the tuned \FS{} finishes everything on chains, diamonds, strongly
connected components, empty graphs, precedence forests, dense and sparse DAGs,
and the historical precedence-knapsack set; the parametric backend matches or
beats it on time in every one of those families --- a tie on the trivial ones,
a factor of $6$ on the historical instances and a factor of $29$ on the sparse
DAGs. It solves $3$ of $9$ layered
DAGs where the parametric backend solves all nine and Dinkelbach eight. On
disconnected graphs and on the MineLib-derived instances it solves nothing
within the limit, and neither does Dinkelbach; only the parametric backend does.
The claim that the \FS{} and Dinkelbach are comparable on the applied instances,
which an earlier reading of a handful of cases suggested, is not supported by
the frozen test set: on the historical precedence-knapsack family both finish
everything, and the medians are $0.087$\,s for Dinkelbach against $0.655$\,s.

One further correction belongs here rather than in a footnote. The ablation
shows two settings that the parameter search rejected and that are faster than
the frozen configuration on the paired comparison: the sink initial basis
($3.33$ against $2.83$) and transitive reduction enabled ($3.46$). The search
was run and frozen honestly, on training and validation data, but its choice is
not the best one available in the space it searched, and a second search with a
larger budget is the obvious next step.

\subsection{What changed, and why it matters more than the parameters}

The first measurements of this project had the \FS{} implementation one to three
orders of magnitude behind both baselines, with time-outs on medium instances.
The gap was not closed by tuning the pricing rule, which barely moves the
result, but by three structural corrections identified from the degeneracy
analysis of section~\ref{sec:robustness}: releasing the anti-cycling fallback
instead of letting it become the operating policy, restricting the ratio-test
scan to the arcs that can actually block, and stopping that scan at the first
zero-step blocker. The last two reproduce techniques already present in the
group's historical prototype, which had no anti-cycling rule at all and carried
an explicit feasibility shortcut for exactly this purpose.

The lesson recorded here is that in this problem class the entering rule is
close to irrelevant while the treatment of degeneracy dominates. Parameter
search over pricing schemes, block sizes and candidate limits produced
differences within a few per cent; the degeneracy-related settings produced
factors.

\subsection{What the measurements do not support}

They do not license any statement about the state of the art in maximum flow.
The flow backends use a Dinic implementation written for this project; they do
not use pseudoflow, highest-label push-relabel with gap and global relabelling,
or the amortised flow reuse of a proper parametric algorithm. A comparison
against a competitive flow code would very likely widen the gap on the families
where \FS{} already loses, and could also narrow the cases where it currently
wins.

They also do not support any claim of novelty. The audit of the literature has
not been carried out in this phase, by explicit decision, and
\file{docs/claims.md} keeps every novelty row in the \emph{to be verified}
state.

\subsection{What follows}

Two directions are supported by the measurements rather than by preference.
First, pricing and forest rebuilding are the dominant costs of a pivot, together
about four fifths of the solver's time on the hardest family, and both are
recomputed in full where only a small part of the graph changed: incremental
pricing and updates confined to the modified paths are the natural next step.
Second, and more important, the length of the pivot path is what actually
separates this method from the parametric backend, and neither the dual
formulation nor the leaving rules touched it.

A third direction is now ruled out rather than recommended. The dual simplex was
implemented and measured on a dedicated campaign precisely because the primal is
so degenerate, and it does remove the degeneracy it was predicted to remove; but
on the validation split it does not beat the primal, and the encouraging factors
observed on three hard instances did not survive the larger sample. That is the
reason the split exists.

\section{Reproducibility}
\label{sec:reproducibility}

Every artefact in this report is derived from files stored in the project.
\file{tools/report/analyze.py} first checks each campaign manifest for
completeness --- planned runs against stored runs, for the build the manifest
describes --- and refuses to emit tables for a campaign that fails the check. It
then writes the tables and figure data used here, the per-run CSV files in
\file{experiments/analysis/}, and \texttt{provenance.json} recording, for each
campaign, the build fingerprint, the run counts, the status histogram and a hash
of the run records.

A run is identified by instance content hash, task, capacity mode, resolved
configuration, build source hash and replica index. A finished run is never
overwritten and never recomputed: a different build or configuration produces a
different identifier. Interrupting a campaign and restarting it resumes from the
stored state.

The commands that rebuild the whole chain --- environment, instance import and
generation, frozen splits, campaigns, tuning, analysis and this document --- are
listed in \file{docs/reproducibility.md}. A reduced reproduction that repeats
one significant comparison in a few minutes is available as
\file{configs/suites/smoke.json}.

The historical material was read, never modified: archives were expanded into a
temporary directory outside the project, and
\file{provenance/original\_manifest.sha256} records size and SHA-256 of every
original file used, in a form that \texttt{sha256sum --check} accepts.

\appendix
\section{Complete tables and campaign details}
\label{sec:appendix}

\generated{table_pilot}

The pilot campaign fixed the sizes, the time limit and the number of replicas of
the later campaigns. It also exposed a protocol defect that was corrected before
any measurement was kept: the time limit was enforced by the \FS{} pivot loop but
not by the flow backends, so a single run had exceeded the budget. The limit is
now checked before every maximum flow, and the pilot was rerun from scratch with
the corrected build. Since a run's identity includes the build hash, the earlier
results could not be silently reused.

Per-run records, including failed and timed-out runs, are in
\file{experiments/analysis/*.csv}; the campaign manifests in
\file{experiments/runs/*/manifest.json} additionally store the resolved
configuration of every run, its statistics counters and the outcome of its
verification. The completeness report that gates the generation of every table
in this document is \file{experiments/analysis/completeness.json}.


\bibliographystyle{plainnat}
\bibliography{references}
\end{document}
